\chapter{TLS: Securing the Web}

Every chapter so far has built a single component. AES is a permutation on $128$-bit
blocks. GCM turns it into authenticated encryption. Diffie--Hellman lets two
strangers agree on a secret over a monitored channel, and elliptic curves make that
agreement small and fast. Digital signatures prove authorship, and certificates
(Chapter~24) bind a public key to a domain name you can read. TLS is the protocol
that assembles all of these pieces to secure the web.

When your browser shows a padlock and the address bar reads
\texttt{https://}, you are speaking \textbf{TLS}: Transport Layer Security. It is
the protocol underneath HTTPS, but also underneath secure email (SMTP over TLS),
encrypted DNS, most VPNs, the control channel of countless mobile apps, and the
back-end traffic of essentially every cloud service. Trillions of TLS handshakes
happen every day. If you understand one applied protocol completely before you
leave this book, make it this one, because TLS is where the abstractions become
load-bearing: a single mismatched assumption between two of the gears becomes a
real attack with a logo and a CVE number.

We will focus on \textbf{TLS 1.3}, standardized in 2018 as RFC 8446. It is the
version you should deploy, it is dramatically cleaner than its predecessors, and
its design choices read like a catalogue of lessons learned the hard way from two
decades of attacks on TLS 1.0 through 1.2. We will trace a full handshake with
small numbers, watch the key schedule derive a dozen keys from one shared secret,
see how the record layer carries encrypted application data, and understand
exactly which property each message protects. Then we will revisit the famous
disasters --- BEAST, POODLE, Heartbleed, and a few cousins --- because the surest
way to appreciate why TLS 1.3 looks the way it does is to see what the old designs
let an attacker do.

\chapterline

\section{What TLS provides, and against whom}

Before any protocol, pin down the threat model. TLS assumes the worst plausible
network: the \textbf{Dolev--Yao} attacker. Eve is not merely listening. She sits
on the wire and can read, drop, reorder, replay, and inject any packet she likes.
She controls the routers, the coffee-shop access point, possibly the ISP. The
only thing she cannot do is break the cryptographic primitives or read memory
inside Alice's and Bob's machines. This is the realistic model for the public
Internet, and TLS is built to survive it.

Against that attacker, a completed TLS connection between a client (Alice's
browser) and a server (Bob's web server) promises three things.

\begin{itemize}
\item \textbf{Confidentiality.} Everything after the handshake --- URLs, cookies,
form data, page contents --- is encrypted. Eve sees ciphertext and traffic sizes,
not plaintext.
\item \textbf{Integrity.} Eve cannot alter a single byte of the conversation
without detection. Every record is authenticated; a flipped bit makes decryption
fail loudly.
\item \textbf{Authentication.} Alice's browser cryptographically verifies that it
is really talking to the server that owns \texttt{example.com}, not to Eve
impersonating it. This relies on certificates and the PKI from
the previous chapter. (By default only the server is authenticated; client
authentication is optional and used in enterprise and mutual-TLS settings.)
\end{itemize}

\begin{focusbox}
\textbf{What TLS does \emph{not} promise.} It does not hide \emph{that} you talked
to a server, nor how much, nor when. The server's identity leaks through the SNI
field (see later) and through the IP address itself. TLS protects the
\emph{contents and integrity} of a channel to an \emph{authenticated} endpoint. It
says nothing about whether that endpoint is trustworthy --- a phishing site can
get a perfectly valid certificate for \texttt{paypa1.com}. Anonymity is the job of
Tor; endpoint trust is the job of the user and the browser's UI.
\end{focusbox}

There is a fourth property TLS 1.3 takes seriously enough that it deserves a
name of its own: \textbf{forward secrecy}. If, a year from now, Eve breaks into
Bob's server and steals his long-term private key, she should \emph{not} be able
to decrypt the traffic she recorded today. We met this idea in Chapter~23; TLS
1.3 makes it mandatory. We will see exactly how.

\begin{intuitionbox}
Think of a TLS connection as a sealed pneumatic tube between two buildings. The
handshake is the part where the two buildings, shouting across a crowded plaza
full of eavesdroppers and imposters, agree on a tube colour nobody else can guess
and check each other's ID badges. The record protocol is the steady stream of
sealed canisters that follow. Eve can watch canisters fly past and even smash one,
but she cannot open one, forge one, or convince either building she is the other.
\end{intuitionbox}

\section{The shape of a handshake}

A TLS connection has two phases. The \textbf{handshake} runs first: it
establishes shared keys and authenticates the server, using public-key
cryptography. Then the \textbf{record protocol} takes over: it carries the actual
application data under fast symmetric authenticated encryption. The handshake is
the expensive part (a few public-key operations and one network round trip); the
record protocol is cheap and runs for the life of the connection.

The headline achievement of TLS 1.3 is that the handshake takes a single round
trip --- ``1-RTT'' --- before the client can send encrypted application data.
TLS 1.2 needed two. The trick is that the client \emph{guesses} which key-exchange
group the server will accept and sends its key share immediately, in its very
first message. If it guesses right (it almost always does), the server can finish
the key agreement in one reply. Figure~\ref{ch25:fig-handshake-comparison} contrasts the legacy TLS 1.2 two-round-trip flow
with the modern TLS 1.3 one-round-trip protocol. We will then walk each message in detail.

\begin{figure}[H]
\centering
\begin{tikzpicture}[font=\scriptsize, >=stealth]
  % Panel (a): TLS 1.2 (Legacy 2-RTT)
  \node[font=\footnotesize\bfseries, text=structurepurple!90!black] at (3.2, 0.5) {(a) Legacy TLS 1.2 (2-RTT Handshake)};
  
  % Lifelines 1.2
  \node[draw=aliceblue!70, fill=aliceblue!15, rounded corners=2pt, minimum width=1.8cm, minimum height=0.6cm] (c12) at (1.0, 0) {Client};
  \node[draw=bobgreen!70, fill=bobgreen!15, rounded corners=2pt, minimum width=1.8cm, minimum height=0.6cm] (s12) at (5.4, 0) {Server};
  \draw[thick, gray!60] (c12) -- (1.0, -7.2);
  \draw[thick, gray!60] (s12) -- (5.4, -7.2);

  % RTT markers 1.2
  \draw[<->, gray!80, thin] (0.0, -0.6) -- node[left, font=\tiny]{1 RTT} (0.0, -3.2);
  \draw[<->, gray!80, thin] (0.0, -3.5) -- node[left, font=\tiny]{2 RTT} (0.0, -5.9);

  % Flight 1 (ClientHello)
  \draw[->, thick, aliceblue!80!black] (1.0, -0.6) -- node[above=1pt, font=\tiny, sloped]{ClientHello (ciphers, nonces)} (5.4, -1.0);
  
  % Flight 2 (ServerHello, Cert, KEX, Done)
  \draw[->, thick, bobgreen!80!black] (5.4, -1.35) -- node[above=1pt, font=\tiny, sloped]{ServerHello, Cert, KEX} (1.0, -1.75);
  \draw[->, thick, bobgreen!80!black] (5.4, -2.0) -- node[above=1pt, font=\tiny, sloped]{ServerHelloDone} (1.0, -2.4);
  \node[font=\tiny, trapred!90!black, fill=white, inner sep=1pt] at (3.2, -2.85) {[\textit{Certificate sent in cleartext}]};

  % Flight 3 (CKE, CCS, Finished)
  \draw[->, thick, aliceblue!80!black] (1.0, -3.5) -- node[above=1pt, font=\tiny, sloped]{ClientKeyEx, ChangeCipherSpec} (5.4, -3.9);
  \draw[->, thick, aliceblue!80!black] (1.0, -4.15) -- node[above=1pt, font=\tiny, sloped]{Finished (Encrypted)} (5.4, -4.55);

  % Flight 4 (CCS, Finished)
  \draw[->, thick, bobgreen!80!black] (5.4, -4.9) -- node[above=1pt, font=\tiny, sloped]{ChangeCipherSpec, Finished} (1.0, -5.3);

  % Application Data (2-RTT)
  \draw[<->, thick, structurepurple] (1.0, -6.2) -- node[above, font=\tiny\bfseries]{[Encrypted Application Data]} (5.4, -6.2);
  \node[font=\tiny, structurepurple] at (3.2, -6.6) {First app data delivered after \textbf{2 RTTs}};

  % Vertical separator
  \draw[dashed, gray!40, thick] (6.6, 0.5) -- (6.6, -7.2);

  % Panel (b): TLS 1.3 (1-RTT)
  \node[font=\footnotesize\bfseries, text=structurepurple!90!black] at (10.2, 0.5) {(b) Modern TLS 1.3 (1-RTT Handshake)};

  % Lifelines 1.3
  \node[draw=aliceblue!70, fill=aliceblue!15, rounded corners=2pt, minimum width=1.8cm, minimum height=0.6cm] (c13) at (7.8, 0) {Client};
  \node[draw=bobgreen!70, fill=bobgreen!15, rounded corners=2pt, minimum width=1.8cm, minimum height=0.6cm] (s13) at (12.6, 0) {Server};
  \draw[thick, gray!60] (c13) -- (7.8, -7.2);
  \draw[thick, gray!60] (s13) -- (12.6, -7.2);

  % RTT markers 1.3
  \draw[<->, gray!80, thin] (7.0, -0.6) -- node[left, font=\tiny]{1 RTT} (7.0, -3.5);

  % Flight 1 (ClientHello + key_share)
  \draw[->, thick, aliceblue!80!black] (7.8, -0.6) -- node[above=1pt, font=\tiny, sloped]{ClientHello + \textbf{key\_share} ($g^a$)} (12.6, -1.0);
  \node[font=\tiny, gray!70, fill=white, inner sep=1pt] at (10.2, -1.3) {Client speculatively offers DH share};

  % Flight 2 (ServerHello + Encrypted Handshake)
  \draw[->, thick, bobgreen!80!black] (12.6, -1.65) -- node[above=1pt, font=\tiny, sloped]{ServerHello + \textbf{key\_share} ($g^b$)} (7.8, -2.05);
  \node[font=\tiny, targetgreen!80!black, fill=white, inner sep=1pt] at (10.2, -2.35) {$\Longrightarrow$ Both sides compute $Z = g^{ab}$};

  \draw[->, thick, structurepurple] (12.6, -2.85) -- node[above=1pt, font=\tiny, sloped]{\{\textbf{EncryptedExtensions}, \textbf{Certificate}\}} (7.8, -3.25);
  \draw[->, thick, structurepurple] (12.6, -3.55) -- node[above=1pt, font=\tiny, sloped]{\{\textbf{CertificateVerify}, \textbf{Finished}\}} (7.8, -3.95);
  \node[font=\tiny, structurepurple, fill=white, inner sep=1pt] at (10.2, -4.25) {[\textit{Certificate is fully encrypted!}]};

  % Flight 3 (Client Finished + App Data)
  \draw[->, thick, structurepurple] (7.8, -4.7) -- node[above=1pt, font=\tiny, sloped]{\{\textbf{Finished}\}} (12.6, -5.1);

  % Application Data (1-RTT)
  \draw[<->, thick, structurepurple] (7.8, -5.8) -- node[above, font=\tiny\bfseries, fill=white, inner sep=1pt]{[Encrypted Application Data]} (12.6, -5.8);
  \node[font=\tiny, targetgreen!80!black] at (10.2, -6.3) {First app data delivered after \textbf{1 RTT}};
  \node[font=\tiny, gray!70] at (10.2, -6.7) {\{...\} = Handshake encrypted};
  \node[font=\tiny, gray!70] at (10.2, -7.0) {[...] = App encrypted};

\end{tikzpicture}
\caption{Comparison of TLS handshakes. In legacy TLS 1.2 (a), public-key parameters and certificates
are negotiated sequentially in cleartext over two round trips before application data can flow.
In TLS 1.3 (b), the client speculatively transmits its Diffie--Hellman key share in the first packet,
collapsing connection latency to 1 RTT while encrypting the server's certificate and extensions.}
\label{ch25:fig-handshake-comparison}
\end{figure}

The notation matters. Messages in the clear are sent before any key exists.
Messages in braces $\{\cdots\}$ are encrypted under \emph{handshake} keys --- a key
the two parties can compute the instant the ServerHello arrives, even though the
server has not yet proven who it is. Messages in brackets $[\cdots]$ ride under the
final \emph{application} keys. That TLS 1.3 encrypts most of the handshake itself
(the certificate, the extensions) is a privacy win over 1.2, where the server's
certificate flew in plaintext for any eavesdropper to log.

\section{Step 1: key exchange (the ClientHello)}

The client opens with a \texttt{ClientHello}. Stripped to essentials it carries:

\begin{itemize}
\item a random nonce $r_C$ (32 bytes), to make every handshake unique;
\item a list of \emph{cipher suites} it supports (e.g.
\texttt{TLS\_AES\_128\_GCM\_SHA256});
\item a list of \emph{supported groups} for key exchange (e.g. X25519,
\texttt{secp256r1}, or finite-field groups);
\item one or more \emph{key shares}: for each group it is betting on, an actual
ephemeral public value. For an elliptic-curve group this is a point; for a
classic finite-field group it is $g^a \bmod p$.
\item the \emph{server name} it wants to reach (the SNI extension), so a server
hosting many sites knows which certificate to present.
\end{itemize}

The key share is the heart of it. The client picks a fresh random secret $a$,
computes its public share, and ships it immediately. Because $a$ is fresh and
thrown away when the connection ends, this is \textbf{ephemeral} Diffie--Hellman ---
the source of forward secrecy. TLS 1.3 \emph{only} allows ephemeral
(EC)DHE key exchange; the old static-RSA key transport, where the client encrypted
a secret under the server's long-term RSA key, is gone, precisely because it had no
forward secrecy.

\subsection{A worked exchange with tiny numbers}

Real TLS uses X25519 or 256-bit curves, with numbers hundreds of bits long. To
\emph{see} the arithmetic, shrink it to a finite-field group you can do by hand.
Take the prime $p = 23$ and generator $g = 5$ (a generator of $\Zns{23}$).

The client picks secret $a = 6$ and sends
\[
A = g^{a} \bmod p = 5^{6} \bmod 23.
\]
Compute $5^2 = 25 \equiv 2$, so $5^4 \equiv 2^2 = 4$, and
$5^6 = 5^4 \cdot 5^2 \equiv 4 \cdot 2 = 8$. The client sends $A = 8$.

The server picks secret $b = 15$ and sends
\[
B = g^{b} \bmod p = 5^{15} \bmod 23.
\]
Build it up: $5^4 \equiv 4$, $5^8 \equiv 4^2 = 16$, $5^{15} = 5^{8} \cdot 5^{4}
\cdot 5^{2} \cdot 5^{1} \equiv 16 \cdot 4 \cdot 2 \cdot 5$. Now $16\cdot4 = 64
\equiv 64 - 2\cdot23 = 18$; $18 \cdot 2 = 36 \equiv 13$; $13 \cdot 5 = 65 \equiv 65
- 2\cdot23 = 19$. So $B = 19$.

Each side now raises the other's share to its own secret:
\[
\text{client: } B^{a} = 19^{6} \bmod 23, \qquad
\text{server: } A^{b} = 8^{15} \bmod 23.
\]
The client: $19 \equiv -4$, so $19^2 \equiv 16$, $19^4 \equiv 16^2 = 256 \equiv
256 - 11\cdot23 = 3$, $19^6 = 19^4\cdot19^2 \equiv 3\cdot16 = 48 \equiv 2$. The
server: $8^2 = 64 \equiv 18$, $8^4 \equiv 18^2 = 324 \equiv 324 - 14\cdot23 = 2$,
$8^8 \equiv 2^2 = 4$, $8^{15} = 8^8\cdot8^4\cdot8^2\cdot8 \equiv 4\cdot2\cdot18\cdot8$.
Then $4\cdot2 = 8$; $8\cdot18 = 144 \equiv 144 - 6\cdot23 = 6$; $6\cdot8 = 48 \equiv
2$. Both get the same shared secret
\[
Z = g^{ab} \bmod p = 2.
\]
Eve, who saw $p=23$, $g=5$, $A=8$, $B=19$, must solve the discrete log to recover
$a$ or $b$. With $p=23$ she can; with a 256-bit curve she cannot, and that gap is
the whole security of the exchange.

\begin{connectionbox}
The cost of computing $g^a \bmod p$ is the modular-exponentiation cost from the
algorithms companion: $O(\log a)$ multiplications by square-and-multiply, which is
why we built up $5^{15}$ from $5^8, 5^4, 5^2$ rather than multiplying fifteen
times. The security rests on the discrete-log assumption (Chapter~18)
or, for X25519, the elliptic-curve discrete-log problem (Chapter~21). The
shared secret $Z$ is \emph{not} used directly as a key --- raw $g^{ab}$ has algebraic
structure and is not uniform. It is run through a key-derivation function, the
subject of Section~\ref{ch25:schedule}, exactly as the ``hash the DH output''
advice in Chapter~23 demanded.
\end{connectionbox}

\section{Step 2: the ServerHello and the handshake key}

The server receives the ClientHello, picks one of the offered groups, generates
its own ephemeral secret $b$, and replies with a \texttt{ServerHello} containing
its own random nonce $r_S$, the chosen cipher suite, and its key share $B$. The
instant this message lands, both parties hold the same shared secret $Z = g^{ab}$
and the same transcript of the two hellos.

From here, \emph{everything} is encrypted. The server immediately derives the
handshake traffic keys (next section) and sends the rest of its messages under
them. This is a sharp break from TLS 1.2, where the certificate was plaintext.

\begin{trapbox}
A subtle point that trips people up: at this moment the client has a confidential,
integrity-protected channel to \emph{somebody}, but it does not yet know
\emph{who}. The Diffie--Hellman exchange authenticates nothing; an active Eve could
have run her own DH with the client. Confidentiality without authentication is a
trap --- it gives a secure channel to the attacker. Authentication is the job of
the next two messages (Certificate, CertificateVerify). Until those verify, the
client must treat the peer as unproven.
\end{trapbox}

\section{Step 3: the key schedule}
\label{ch25:schedule}

We have a shared secret $Z$. We need many keys: one to encrypt the server's
handshake messages, one for the client's, one for application data in each
direction, plus keys to compute the Finished MACs. TLS 1.3 manufactures all of
them from $Z$ with a single tool: \textbf{HKDF}, the HMAC-based key-derivation
function from Chapter~23 (built on HMAC from Chapter~15).

HKDF has two halves. \emph{Extract} takes a salt and some input keying material
and squeezes out a uniform pseudorandom key:
\[
\mathrm{PRK} = \mathrm{HKDF\text{-}Extract}(\text{salt}, \text{IKM})
= \operatorname{HMAC}(\text{salt}, \text{IKM}).
\]
\emph{Expand} stretches that key into as many labelled output keys as you want,
binding each to a label and to the handshake transcript:
\[
\mathrm{HKDF\text{-}Expand\text{-}Label}(\text{secret}, \text{label}, \text{context}, L).
\]
The label (a short ASCII string like \texttt{"c hs traffic"}) guarantees that two
keys derived for different purposes are independent even though they came from the
same secret. The context is usually a hash of the handshake transcript so far, so
the keys are cryptographically bound to \emph{exactly} the messages that have been
exchanged.

TLS 1.3 arranges these derivations in a \textbf{schedule}: a chain of Extract and
Expand steps, each feeding the next (Figure~\ref{ch25:fig-keyschedule}).

\begin{figure}[H]
\centering
\begin{tikzpicture}[font=\scriptsize, >=stealth, node distance=10mm]
  % Stage 1: Early Secret
  \node[draw=quotientorange!80, fill=quotientorange!15, rounded corners=2pt, font=\tiny\bfseries, align=center, text width=2.4cm] (psk) at (-4.2, 3.5) {Input Key Material:\\$\text{PSK}$ or $0$};
  \node[draw=structurepurple!80, fill=structurepurple!15, rounded corners=2pt, font=\scriptsize\bfseries, align=center, text width=3.8cm] (es) at (0, 3.5) {
    \textbf{Early Secret}\\
    $\operatorname{HKDF-Extract}(0, \text{PSK})$
  };
  \draw[->, thick, quotientorange!80!black] (psk) -- (es);

  % Outputs from Early Secret
  \node[draw=aliceblue!80, fill=aliceblue!15, rounded corners=2pt, font=\tiny, align=left, text width=4.5cm] (early_out) at (5.55, 3.5) {
    $\bullet$ \texttt{client\_early\_traffic\_secret} (0-RTT)\\
    $\bullet$ \texttt{early\_exporter\_master\_secret}
  };
  \draw[->, thick, gray!70] (es) -- node[above=2pt, font=\tiny]{Expand} (early_out);

  % Derivation link 1 -> 2
  \node[draw=gray!50, fill=gray!10, rounded corners=1pt, font=\tiny] (d1) at (0, 2.2) {$\operatorname{Derive-Secret}(\text{ES}, \text{"derived"})$};
  \draw[->, thick, structurepurple!80] (es) -- (d1);

  % Stage 2: Handshake Secret
  \node[draw=quotientorange!80, fill=quotientorange!15, rounded corners=2pt, font=\tiny\bfseries, align=center, text width=2.4cm] (dhe) at (-4.2, 0.9) {ECDHE Secret:\\$Z = g^{ab} = a b G$};
  \node[draw=structurepurple!80, fill=structurepurple!15, rounded corners=2pt, font=\scriptsize\bfseries, align=center, text width=3.8cm] (hs) at (0, 0.9) {
    \textbf{Handshake Secret}\\
    $\operatorname{HKDF-Extract}(\text{Salt}_1, Z)$
  };
  \draw[->, thick, structurepurple!80] (d1) -- (hs);
  \draw[->, thick, quotientorange!80!black] (dhe) -- (hs);

  % Outputs from Handshake Secret
  \node[draw=targetgreen!80, fill=targetgreen!10, rounded corners=2pt, font=\tiny, align=left, text width=4.5cm] (hs_out) at (5.55, 0.9) {
    $\bullet$ \texttt{client\_handshake\_traffic\_secret}\\
    $\bullet$ \texttt{server\_handshake\_traffic\_secret}\\
    {\color{gray!80}$\to$ Protects EncryptedExtensions, Cert, Finished}
  };
  \draw[->, thick, gray!70] (hs) -- node[above=2pt, font=\tiny]{Expand} (hs_out);

  % Derivation link 2 -> 3
  \node[draw=gray!50, fill=gray!10, rounded corners=1pt, font=\tiny] (d2) at (0, -0.4) {$\operatorname{Derive-Secret}(\text{HS}, \text{"derived"})$};
  \draw[->, thick, structurepurple!80] (hs) -- (d2);

  % Stage 3: Master Secret
  \node[draw=quotientorange!80, fill=quotientorange!15, rounded corners=2pt, font=\tiny\bfseries, align=center, text width=2.4cm] (zero_ikm) at (-4.2, -1.7) {Empty Input:\\$0$};
  \node[draw=structurepurple!80, fill=structurepurple!15, rounded corners=2pt, font=\scriptsize\bfseries, align=center, text width=3.8cm] (ms) at (0, -1.7) {
    \textbf{Master Secret}\\
    $\operatorname{HKDF-Extract}(\text{Salt}_2, 0)$
  };
  \draw[->, thick, structurepurple!80] (d2) -- (ms);
  \draw[->, thick, quotientorange!80!black] (zero_ikm) -- (ms);

  % Outputs from Master Secret
  \node[draw=structurepurple!80, fill=structurepurple!10, rounded corners=2pt, font=\tiny, align=left, text width=4.5cm] (ms_out) at (5.55, -1.7) {
    $\bullet$ \texttt{client\_application\_traffic\_secret\_0}\\
    $\bullet$ \texttt{server\_application\_traffic\_secret\_0}\\
    $\bullet$ \texttt{exporter\_master\_secret}\\
    $\bullet$ \texttt{resumption\_master\_secret}
  };
  \draw[->, thick, gray!70] (ms) -- node[above=2pt, font=\tiny]{Expand} (ms_out);

  % Context annotation
  \node[draw=gray!50, fill=white, rounded corners=2pt, font=\tiny, align=center, text width=12.5cm] at (0.4, -3.0) {
    \textbf{Transcript Binding}: Each $\operatorname{HKDF-Expand-Label}(\text{Secret}, \text{Label}, \text{Context}, L)$ mixes the hash of all prior handshake messages\\
    $\text{Context} = \Hash(\text{ClientHello} \concat \cdots)$ into the derived key, preventing downgrade attacks and transcript tampering.
  };
\end{tikzpicture}
\caption{The TLS 1.3 HKDF key derivation schedule. Entropy from the optional pre-shared key (PSK)
and ephemeral Diffie--Hellman exchange ($Z$) is absorbed in stages via \texttt{HKDF-Extract}.
Intermediate secrets are expanded into cryptographically isolated traffic keys for the handshake,
application data, and future session resumption.}
\label{ch25:fig-keyschedule}
\end{figure}

The \textbf{Early Secret} is extracted from a pre-shared key (PSK) if one exists,
or from a string of zeros if not. The \textbf{Handshake Secret} mixes in the actual
Diffie--Hellman output $Z$; from it the two handshake traffic keys are expanded,
and the server starts using its key right away. After the handshake completes, the
\textbf{Master Secret} (extracted with the DH input already consumed, so the input
is zeros) yields the application traffic keys that protect real data, plus a
resumption secret for next time.

\begin{whybox}
Why a chain instead of one Expand from $Z$? Three reasons, each a lesson from
older protocols. \textbf{(1) Domain separation:} distinct labels make the handshake
keys provably independent from the application keys, so compromising one phase
does not leak another. \textbf{(2) Transcript binding:} because each key is expanded
with the hash of the messages so far, an attacker who tampers with any handshake
message produces different keys on the two sides, and the Finished check (below)
fails. \textbf{(3) PSK and DH composition:} the Extract steps let TLS fold in a
pre-shared key \emph{and} fresh DH cleanly, so resumed connections still get forward
secrecy from the new $g^{ab}$. TLS 1.2 derived everything in one ad-hoc PRF call
and paid for it with attacks like Triple Handshake that this structure closes.
\end{whybox}

\begin{implementbox}
A toy of the TLS 1.3 key schedule using Python's \texttt{hmac} and
\texttt{hashlib}. It mirrors the real structure but is radically simplified --- no real
labels, no transcript encoding, fixed lengths. \textbf{Learning only: never use
this to secure anything.}

\begin{lstlisting}[language=Python]
import hmac, hashlib

def hkdf_extract(salt, ikm):
    return hmac.new(salt, ikm, hashlib.sha256).digest()

def hkdf_expand_label(secret, label, context, length=32):
    # Real TLS encodes (length, "tls13 "+label, context) into the info field.
    info = b"tls13 " + label + b"|" + context
    out, t, i = b"", b"", 1
    while len(out) < length:
        t = hmac.new(secret, t + info + bytes([i]), hashlib.sha256).digest()
        out += t
        i += 1
    return out[:length]

def derive_secret(secret, label, transcript_hash):
    return hkdf_expand_label(secret, label, transcript_hash)

zeros     = b"\x00" * 32
dh_shared = b"...the g^{ab} bytes go here..."             # from the (EC)DHE step
th = hashlib.sha256(b"ClientHello||ServerHello").digest() # toy transcript hash

early_secret     = hkdf_extract(b"\x00", zeros)           # no PSK -> zeros
handshake_secret = hkdf_extract(
    derive_secret(early_secret, b"derived", b""), dh_shared)
c_hs_traffic     = derive_secret(handshake_secret, b"c hs traffic", th)
s_hs_traffic     = derive_secret(handshake_secret, b"s hs traffic", th)
master_secret    = hkdf_extract(
    derive_secret(handshake_secret, b"derived", b""), zeros)
# Application keys come from master_secret once the transcript is final.
print("client handshake key:", c_hs_traffic.hex())
\end{lstlisting}
\end{implementbox}

\section{Step 4: authentication (Certificate and CertificateVerify)}

Now the server proves it is who the client wanted. It sends, all under the
handshake key:

\begin{itemize}
\item \texttt{EncryptedExtensions} --- assorted negotiated parameters that do not
need to be public.
\item \texttt{Certificate} --- the server's X.509 certificate chain: its own
certificate (binding its public key to its domain name) followed by any
intermediate CA certificates up toward a root the client trusts. This is the PKI
machinery from Chapter~24.
\item \texttt{CertificateVerify} --- the authenticating step. The server takes the entire
handshake transcript hashed so far, signs it with the \emph{private key
corresponding to the public key in its certificate}, and sends the signature.
\end{itemize}

The CertificateVerify message is what fuses identity to the live connection. The
certificate alone proves nothing: anyone can copy Bob's certificate, since it is
public. What an attacker cannot do is produce a valid signature over
\emph{this particular transcript} --- which includes the client's fresh random and
the ephemeral key shares --- without Bob's private key. The signature is fresh,
bound to this handshake, and unforgeable.

\begin{whybox}
\textbf{Why signing the transcript authenticates the key exchange.} The transcript
the server signs contains the ServerHello, hence the server's key share $B$, hence
(together with the client's $A$) the shared secret $Z$. So the signature
simultaneously says ``I am the holder of the certified private key'' \emph{and}
``the key share you are about to derive your keys from is mine, not an
attacker's.'' This is the textbook \emph{signed Diffie--Hellman} pattern from
Chapter~23: ephemeral DH gives forward secrecy, a signature over the DH
values gives authentication, and binding them in one transcript stops a
man-in-the-middle from splicing her own key share in. If active Eve had injected
her share, the transcript Bob signs would not match the one Alice computed, and
verification --- or the Finished check --- fails.
\end{whybox}

The client verifies the certificate chain (validity dates, name match against the
SNI it requested, a path to a trusted root, revocation status), then verifies the
CertificateVerify signature against the public key in that certificate. Only if
both pass does the client believe it is talking to the real server.

\section{Step 5: Finished (the MAC that seals the handshake)}

The server sends one more handshake message: \texttt{Finished}. It is a MAC ---
specifically $\operatorname{HMAC}$ under a key derived from the handshake secret --- computed over
the \emph{entire transcript up to this point}. The client, holding the same key
and transcript, recomputes the MAC and checks it. Then the client sends its own
\texttt{Finished}, a MAC over the transcript including the server's Finished, and
the server checks that.

These two MACs are the integrity capstone of the handshake. Their value is subtle
and worth stating plainly.

\begin{whybox}
\textbf{Why Finished defeats tampering and downgrades.} Every handshake message ---
the cipher suite list, the supported groups, the key shares, the certificate ---
feeds into the transcript hash. The Finished MAC is computed over that hash with a
key that an attacker cannot derive (it depends on $g^{ab}$). So if Eve altered
\emph{any} earlier message --- swapped a strong cipher suite for a weak one,
stripped a group, flipped a bit in a key share --- the two sides' transcripts
diverge, their Finished MACs do not match, and the handshake aborts. In particular
this is TLS 1.3's \textbf{downgrade protection}: even though the negotiation
happens in the clear, its result is authenticated after the fact. Eve cannot trick
Alice and Bob into using weaker crypto than they both support, because the moment
she edits the negotiation she breaks the Finished check. The lesson was learned the
hard way: TLS 1.2's negotiation was not fully authenticated, and attacks like
FREAK and Logjam forced servers down to deliberately weakened ``export-grade'' keys
precisely by tampering with the unprotected handshake.
\end{whybox}

\begin{focusbox}
\textbf{Putting the pieces together.} Look at what each message bought:
\begin{itemize}
\item ClientHello / ServerHello key shares $\Rightarrow$ a \emph{shared secret}
(confidentiality, forward secrecy).
\item Certificate + CertificateVerify $\Rightarrow$ the server is \emph{who it
claims} (authentication).
\item Finished (both) $\Rightarrow$ \emph{nobody tampered} with the negotiation
(integrity, downgrade protection).
\end{itemize}
Confidentiality, authentication, integrity --- the three promises of the first
section, each pinned to a specific message, each built from a primitive you
already understand.
\end{focusbox}

\section{The record protocol and AEAD}

With the handshake done, the connection switches to the \textbf{record protocol}.
Application data is split into records (up to $2^{14}$ bytes of plaintext each),
and each record is sealed with an \textbf{AEAD} cipher --- authenticated
encryption with associated data --- from Chapter~16. TLS 1.3 permits only
AEAD suites: AES-GCM, AES-CCM, or ChaCha20-Poly1305. No more
encrypt-then-separately-MAC bolt-ons, no unauthenticated CBC; the cipher provides
confidentiality and integrity in one operation.

An AEAD seal takes four inputs: the traffic key $k$, a nonce, the plaintext, and
some associated data that is authenticated but not encrypted. It returns
ciphertext plus an authentication tag:
\[
\mathrm{Seal}_k(\text{nonce}, \text{AD}, \text{plaintext}) \to (\text{ciphertext},
\text{tag}).
\]
Decryption recomputes the tag and refuses to output plaintext if it does not
match. A single flipped bit anywhere in the record makes the tag check fail and
the connection tears down. That is the integrity guarantee, per record.

\subsection{Nonces and the sequence number}

GCM and ChaCha20-Poly1305 are catastrophically insecure if you ever reuse a
(key, nonce) pair --- repeat a nonce under AES-GCM and you leak the XOR of two
plaintexts and, worse, the authentication subkey, letting an attacker forge. TLS
1.3 sidesteps this elegantly. It does not transmit nonces. Instead each side keeps
a 64-bit \textbf{sequence number} that starts at zero and increments by one per
record. The per-record nonce is that counter XORed into a fixed random
\texttt{iv} derived from the traffic secret in the key schedule.

\begin{focusbox}
\textbf{Nonce $=$ static\_iv $\xor$ (sequence number).} Because the sequence number
never repeats within a connection and the \texttt{iv} is unique per direction per
key, no (key, nonce) pair ever recurs. The counter is implicit --- never sent --- so
it also costs zero bytes on the wire. Tiny example: with the last two bytes of the
$\mathtt{static\_iv}$ equal to $\mathtt{A5C3}$ and record numbers $0,1,2$, the
nonces end in $\mathtt{A5C3}, \mathtt{A5C2}, \mathtt{A5C1}$ --- distinct by
construction. When the sequence number would overflow, TLS forces a key update
(re-keying) rather than risk a repeat.
\end{focusbox}

\subsection{What the record header authenticates}

In TLS 1.3 the record's outer header is tiny: a content-type byte (set to
\texttt{application\_data} even for handshake records, to hide what is inside), a
legacy version field, and the length. The real content type is appended to the
plaintext \emph{before} encryption and recovered after decryption, so an
eavesdropper cannot tell handshake records from data records from alerts. The outer
header is fed in as the AEAD associated data, so its length field is authenticated
even though it travels in the clear: an attacker cannot lie about the record's
length without breaking the tag (Figure~\ref{ch25:fig-aead-record}).

\begin{figure}[H]
\centering
\begin{tikzpicture}[font=\scriptsize, >=Stealth]
  % Nonce derivation (left)
  \node[draw=sourceblue!60, fill=sourceblue!12, rounded corners=2pt, align=center, inner sep=3pt, minimum width=2.8cm] (iv) at (0,2.2) {12-byte \texttt{static\_iv}\\(from HKDF)};
  \node[draw=sourceblue!60, fill=sourceblue!12, rounded corners=2pt, align=center, inner sep=3pt, minimum width=2.8cm] (seq) at (0,1.0) {64-bit seq number\\(implicit counter)};
  \node[draw=structurepurple!80, fill=structurepurple!15, circle, inner sep=1pt] (xor) at (2.2,1.6) {$\xor$};
  \node[draw=structurepurple!70, fill=structurepurple!12, rounded corners=2pt, align=center, inner sep=3pt] (nonce) at (4.0,1.6) {Per-record\\12B Nonce};

  \draw[->, thick, gray!70] (iv.east) -| (xor.north);
  \draw[->, thick, gray!70] (seq.east) -| (xor.south);
  \draw[->, thick, structurepurple!80!black] (xor.east) -- (nonce.west);

  % Inner plaintext (top right)
  \node[draw=bobgreen!60, fill=bobgreen!12, rounded corners=2pt, align=center, inner sep=3pt, minimum width=2.8cm] (pt) at (7.5,2.2) {Plaintext payload};
  \node[draw=bobgreen!60, fill=bobgreen!20, rounded corners=2pt, align=center, inner sep=3pt, minimum width=1.8cm] (ct_type) at (10.0,2.2) {Inner Type\\(\texttt{0x16}/\texttt{0x17})};

  % AEAD Seal box
  \node[draw=quotientorange!80, fill=quotientorange!15, rounded corners=2pt, align=center, inner sep=4pt, minimum width=3.8cm] (aead) at (8.5,1.0) {\textbf{AEAD Seal}$_k(\text{Nonce}, \text{AD}, \text{Plaintext})$};

  \draw[->, thick, structurepurple!80!black] (nonce.east) -- (aead.west);
  \draw[->, thick, bobgreen!80!black] (pt.south) -- (aead.north);
  \draw[->, thick, bobgreen!80!black] (ct_type.south) -- (aead.north east);

  % Wire record (bottom)
  \node[draw=gray!70, fill=gray!12, rounded corners=2pt, align=center, inner sep=4pt, minimum width=3.2cm] (hdr) at (3.6,-0.4) {\textbf{Outer Header (AD)}\\\texttt{type=0x17}, \texttt{ver=0x0303}, \texttt{len}};
  \node[draw=targetgreen!70, fill=targetgreen!12, rounded corners=2pt, align=center, inner sep=4pt, minimum width=4.5cm] (ciph) at (7.6,-0.4) {\textbf{Encrypted Inner Record}\\(Ciphertext)};
  \node[draw=trapred!70, fill=trapred!12, rounded corners=2pt, align=center, inner sep=4pt, minimum width=2.3cm] (tag) at (11.2,-0.4) {\textbf{16B Auth Tag}\\(MAC over all)};

  \draw[->, thick, gray!70, dashed] (hdr.north) |- node[pos=0.25, left, font=\tiny]{Authenticated (AD)} (aead.south west);
  \draw[->, thick, quotientorange!80!black] (aead.south) -- (ciph.north);
  \draw[->, thick, quotientorange!80!black] (aead.south east) -- (tag.north);
\end{tikzpicture}
\caption{TLS 1.3 record-layer AEAD sealing. The 64-bit sequence number is XORed with the static IV to form a unique 12-byte nonce without transmitting a counter on the wire. The true content type is hidden inside the encrypted payload, while the cleartext outer header is authenticated as Associated Data (AD).}
\label{ch25:fig-aead-record}
\end{figure}

\begin{applicationbox}
\textbf{This is the traffic under your padlock.} When you load a page over HTTPS,
the TLS handshake runs once (often resumed in well under a millisecond on repeat
visits), and then every HTTP request and response --- including HTTP/2 and HTTP/3
framing --- rides inside TLS records as AEAD-sealed canisters. On a modern CPU
with AES-NI, AES-GCM runs at gigabytes per second, so the record layer is
essentially free; the handshake's public-key math is the only part you can feel,
and even that is one round trip. ChaCha20-Poly1305 exists for the same reason but
on phones and embedded chips without AES hardware, where its software speed and
constant-time design (no secret-dependent table lookups) win. Google deployed it
across Android and Chrome for exactly that mobile-battery, side-channel-safety
combination.
\end{applicationbox}

\section{Forward secrecy, made concrete}

We have invoked forward secrecy several times; pin it down. A protocol has forward
secrecy if compromise of long-term secrets in the future does not compromise past
session keys. In TLS 1.3 the long-term secret is the server's certificate private
key. The session keys come from the \emph{ephemeral} Diffie--Hellman secret $Z =
g^{ab}$, where $a$ and $b$ are random, used once, and erased when the connection
closes.

\begin{intuitionbox}
Suppose Eve records gigabytes of your encrypted TLS traffic today and files it
away. Five years from now she steals Bob's server, including its certificate
private key. Can she decrypt the old recordings? \textbf{No.} The certificate key
was used only to \emph{sign} the handshake transcript, never to derive any
encryption key. The encryption keys came from $g^{ab}$, and $a$ and $b$ were thrown
away the moment those connections ended. To get $Z$ she would need $a$ or $b$, and
the only thing on the wire was $g^a$ and $g^b$ --- recovering the exponents is the
discrete-log problem she cannot solve. The stolen signing key lets her impersonate
Bob \emph{going forward}, but it cannot retroactively open a single old session.
\end{intuitionbox}

\begin{historybox}
This is the property that static-RSA TLS lacked. TLS grew out of Netscape's Secure
Sockets Layer (SSL~2.0 in 1995 and SSL~3.0 in 1996, designed by Paul Kocher, Taher
ElGamal, and colleagues) and was standardized by the IETF as TLS~1.0 in 1999. In TLS
up to 1.2, a common mode was \emph{RSA key transport}: the client picked the session
secret and encrypted it under the server's long-term RSA public key from the
certificate, making the session key recoverable by anyone holding that RSA private
key --- \emph{forever}. The 2013 Snowden disclosures made the risk concrete: an
adversary who recorded traffic in bulk and later obtained a server's private key (by
subpoena, theft, or coercion) could decrypt years of archived sessions retroactively.
Under editor Eric Rescorla, TLS~1.3 (2018, RFC~8446) finished the push toward
ephemeral (EC)DHE by \emph{removing static RSA key transport from the protocol
entirely}, over objections from enterprises whose monitoring middleboxes relied on
passive decryption with a copy of the server's key. Forward secrecy is no longer an
option you can forget to enable; it is the only thing on offer.
\end{historybox}

\section{Resumption and 0-RTT (and its one sharp edge)}

Re-running a full handshake (with its signature and certificate verification) for
every connection to a site you just visited is wasteful. TLS 1.3 lets the server
issue a \textbf{session ticket} after a handshake: a pre-shared key (PSK) the client
can present next time to skip the public-key authentication. A resumed handshake
is faster and still mixes in fresh DH for forward secrecy (the Handshake Secret
extracts both the PSK and the new $g^{ab}$).

TLS 1.3 also offers \textbf{0-RTT}: using a resumption PSK, the client can send
encrypted application data \emph{in its very first flight}, before any round trip,
shaving the last latency off a repeat visit. It is a genuine speedup, but it comes
with a caveat severe enough that it must be stated loudly.

\begin{trapbox}
\textbf{0-RTT data has no replay protection.} The early data is encrypted under a
key derived from the PSK alone, with no fresh server contribution yet. Eve cannot
read it, but she \emph{can} record the 0-RTT flight and replay it later --- the
server has no per-connection nonce from itself to reject the duplicate. If that
early request is a non-idempotent action (``transfer \$100,'' ``post this
comment''), a replay executes it twice. The rule: \emph{only ever put idempotent,
safe requests in 0-RTT} (a GET, a cacheable read), never a state-changing one.
Application servers must enforce this; many simply forbid 0-RTT for anything but
GETs. It is a textbook case of a performance feature trading away a security
property, and you have to know which.
\end{trapbox}

Figure~\ref{ch25:fig-0rtt-replay} diagrams why this replay vulnerability occurs in practice.

\begin{figure}[H]
\centering
\begin{tikzpicture}[font=\scriptsize, >=stealth]
  % Parties
  \node[draw=aliceblue!80, fill=aliceblue!15, rounded corners=2pt, minimum width=2.2cm, minimum height=0.7cm] (alice) at (0, 0) {Client (Alice)};
  \node[draw=trapred!80, fill=trapred!15, rounded corners=2pt, minimum width=2.2cm, minimum height=0.7cm] (eve) at (5.2, 0) {Attacker (Eve)};
  \node[draw=bobgreen!80, fill=bobgreen!15, rounded corners=2pt, minimum width=2.2cm, minimum height=0.7cm] (server) at (10.4, 0) {Server (Bank)};

  \draw[thick, gray!50] (alice) -- (0, -4.5);
  \draw[thick, gray!50] (eve) -- (5.2, -4.5);
  \draw[thick, gray!50] (server) -- (10.4, -4.5);

  % Legitimate 0-RTT flight
  \draw[->, thick, aliceblue!80!black] (0, -0.85) -- node[above=1pt, font=\tiny, align=center]{1. 0-RTT Flight: ClientHello + PSK\\\{\texttt{POST /transfer?to=Bob\&amt=100}\}} (5.2, -0.85);
  \node[draw=trapred!60, fill=trapred!10, rounded corners=1pt, font=\tiny, align=center] at (5.2, -1.35) {Eve records ciphertext};
  \draw[->, thick, gray!80] (5.2, -1.75) -- (10.4, -1.75);
  \node[anchor=west, font=\tiny, bobgreen!80!black] at (10.55, -1.75) {Transfer \#1};

  % Server reply
  \draw[->, thick, bobgreen!80!black] (10.4, -2.3) -- node[above=1pt, font=\tiny]{2. ServerHello + Finished + HTTP 200 OK} (0, -2.3);

  % Replay flight by Eve
  \draw[->, thick, trapred!90!black, dashed] (5.2, -3.1) -- node[above=1pt, font=\tiny\bfseries, text=trapred!90!black]{3. Eve REPLAYS recorded 0-RTT flight} (10.4, -3.1);
  
  \node[draw=trapred!70, fill=white, rounded corners=2pt, font=\tiny, align=left, text width=4.8cm, inner sep=2pt] at (7.8, -3.95) {
    \textbf{Server Flaw}: Without a fresh server nonce in 0-RTT, server decrypts flight with valid PSK and \textbf{executes transfer \#2}!
  };

  % Rule box
  \node[draw=targetgreen!70, fill=targetgreen!10, rounded corners=2pt, font=\tiny, align=center, text width=12.0cm] at (5.2, -4.8) {
    \textbf{The 0-RTT Security Boundary}: 0-RTT provides \emph{confidentiality} but \textbf{no replay protection}.\\
    Protocols must restrict 0-RTT exclusively to idempotent operations (e.g., HTTP GET read requests), never financial or state-modifying actions.
  };
\end{tikzpicture}
\caption{The 0-RTT early data replay vulnerability in TLS 1.3. Because 0-RTT data is encrypted before
receiving a fresh random nonce from the server, an active eavesdropper can capture and retransmit the flight.
Unless application servers enforce idempotency or anti-replay caches, state-altering requests are duplicated.}
\label{ch25:fig-0rtt-replay}
\end{figure}

\section{Famous TLS attacks}

The streamlined architecture of TLS 1.3 was shaped directly by two decades of hard-won security lessons. Each
of the following attacks broke a real deployment, and each maps to a concrete design
decision in 1.3 that eliminated the vulnerability class. Studying them develops
a sharp eye for where protocol abstractions leak.

\begin{attackbox}
\textbf{BEAST (2011) --- a predictable IV in CBC mode.} TLS 1.0 used CBC mode with a
fatal shortcut: the IV for each record was the \emph{last ciphertext block of the
previous record}, so it was predictable to an attacker watching the wire. Recall
CBC encrypts block $i$ as $C_i = \Enc_k(P_i \xor C_{i-1})$. If Eve can predict
$C_{i-1}$ \emph{and} influence the plaintext $P_i$ (she can, via malicious
JavaScript making the browser send chosen data alongside a secret cookie), she can
test guesses of a secret byte: choose $P_i$ so that $P_i \xor C_{i-1}$ equals a
value she wants to probe, and watch whether the resulting ciphertext block collides
with one she has seen for a known plaintext. Byte by byte she recovers the cookie.
\emph{Fix:} TLS 1.1 made the IV random and unpredictable per record; TLS 1.3
abandons CBC entirely for AEAD, where this whole class of chosen-plaintext games
against the mode does not exist.
\end{attackbox}

\begin{figure}[H]
\centering
\begin{tikzpicture}[font=\small,>=Stealth,node distance=7mm]
  \node (P0) at (0,1.6) {$P_{i-1}$};
  \node (P1) at (3.6,1.6) {$P_i$};
  \node[draw,circle,inner sep=1pt] (X0) at (0,0.8) {$\xor$};
  \node[draw,circle,inner sep=1pt] (X1) at (3.6,0.8) {$\xor$};
  \node[draw,rounded corners,fill=structurepurple!12] (E0) at (0,0) {$\Enc_k$};
  \node[draw,rounded corners,fill=structurepurple!12] (E1) at (3.6,0) {$\Enc_k$};
  \node (C0) at (0,-0.9) {$C_{i-1}$};
  \node (C1) at (3.6,-0.9) {$C_i$};
  \node[evered,font=\scriptsize,align=center] (IV) at (-2.4,0.8) {prev.\ block\\(predictable IV)};
  \draw[->] (P0) -- (X0); \draw[->] (P1) -- (X1);
  \draw[->] (X0) -- (E0); \draw[->] (X1) -- (E1);
  \draw[->] (E0) -- (C0); \draw[->] (E1) -- (C1);
  \draw[->] (IV) -- (X0);
  \draw[->, evered] (C0.south) -- ++(0,-0.35) -- ++(1.8,0) |- (X1.west);
  \node[evered,font=\scriptsize,below=1pt, fill=white, inner sep=1pt] at (1.0,-1.25) {$C_{i-1}$ feeds $C_i$};
\end{tikzpicture}
\caption{CBC chaining as BEAST exploits it: each ciphertext block is the IV for
the next, so an attacker who has seen $C_{i-1}$ knows the value XORed into the
chosen $P_i$ and can probe a secret byte through ciphertext collisions.}
\label{ch25:beast-cbc}
\end{figure}

Concretely (Figure~\ref{ch25:beast-cbc}): because Eve knows $C_{i-1}$ and chooses
$P_i$, she controls $P_i \xor C_{i-1}$ exactly. A random IV unknown to her before
she commits $P_i$ destroys that control, which is why TLS 1.1 patched it and TLS
1.3 dropped the mode.

\begin{attackbox}
\textbf{POODLE (2014) --- padding you cannot trust.} SSL 3.0's CBC padding was
\emph{unauthenticated} and its contents unspecified except for the last byte. Under
MAC-then-encrypt, the receiver decrypts, strips padding, then checks the MAC --- so
the padding is processed before authentication. Eve, as an active man-in-the-middle
running a padding-oracle attack, replaces the last ciphertext block with an earlier
block she wants to decrypt; with probability $1/256$ the mangled padding happens to
be accepted, and that acceptance leaks one plaintext byte. By forcing many
connections (again via injected JavaScript) she peels off a secret byte at a time.
The deeper flaw was \textbf{MAC-then-encrypt}: integrity was checked too late, after
the attacker-controlled padding had already influenced behaviour. \emph{Fix:} stop
using SSL 3.0 (POODLE is also why downgrade protection matters --- the attack often
began by forcing a downgrade to SSL 3.0), and structurally, move to AEAD and to
encrypt-then-MAC ordering, both of which TLS 1.3 mandates. The padding-oracle
lineage runs through Vaudenay's 2002 attack, Lucky Thirteen's timing variant, and
POODLE --- all the same disease.
\end{attackbox}

\begin{attackbox}
\textbf{Heartbleed (2014) --- not a protocol flaw at all.} Heartbleed was a buffer
over-read in OpenSSL's implementation of the TLS heartbeat extension. The heartbeat
request said ``echo back these $N$ bytes,'' and the code trusted $N$ without
checking it against the actual payload length. An attacker sent a heartbeat
claiming a 1-byte payload but a length of $64$KB; the server dutifully copied $64$KB
of \emph{adjacent process memory} into the reply --- which could contain other
users' requests, session cookies, or, devastatingly, the server's private key. No
cryptography was broken; the math was perfect. A single missing bounds check in C
exposed the long-term private keys of a large fraction of the web overnight,
forcing mass certificate revocation and re-issuance. \emph{Lesson:} the protocol is
only as strong as the code that implements it. This is the entire motivation for
Chapters~37--39 --- memory safety, bounds checks, and
minimizing the attack surface are cryptographic concerns even when the
cryptography is flawless. (It is also a sharp reminder of why forward secrecy
matters: servers that used ephemeral key exchange at least did not retroactively
expose past sessions when their keys leaked.)
\end{attackbox}

\begin{attackbox}
\textbf{The downgrade family --- FREAK and Logjam (2015).} For 1990s export
regulations, TLS retained support for deliberately weak ``export-grade'' RSA
($512$-bit) and DH ($512$-bit) parameters. FREAK and Logjam were active attacks
that tampered with the (unauthenticated) negotiation in TLS 1.2 to coax a client
and server into using these weak parameters, after which the attacker factored the
$512$-bit RSA modulus or solved the $512$-bit discrete log offline --- feasible with
modest compute, especially since many servers shared the \emph{same} DH group, so
one big precomputation cracked many sites (Logjam). \emph{Fix:} TLS 1.3 removes
export ciphers and all the legacy weak groups, fixes a short list of strong groups,
and --- the structural cure --- authenticates the entire negotiation via the Finished
MAC, so an attacker who edits the cipher/group choice is caught.
\end{attackbox}

\begin{focusbox}
\textbf{The pattern across all of them.} Predictable IVs (BEAST), authenticate-too-late
padding (POODLE, Lucky Thirteen), unauthenticated negotiation (FREAK, Logjam), and
implementation memory bugs (Heartbleed). TLS 1.3's answers: random per-record
nonces from a counter, AEAD only (encrypt-and-authenticate together), a fully
transcript-bound Finished MAC, and a smaller, simpler surface with fewer options to
get wrong. Almost every line of the 1.3 spec is a ``never again'' to one of these.
\end{focusbox}

\section{A bird's-eye view of the gears}

\begin{connectionbox}
TLS is the integration test for the whole book. The (EC)DHE key exchange is
Diffie--Hellman over the elliptic curves of Chapter~21, giving the shared
secret and forward secrecy. The key schedule is HKDF (Chapter~23), built from the HMAC of
Chapter~15. The server's CertificateVerify is a digital signature (ECDSA or
RSA-PSS or Ed25519) from Chapter~22, verified against a certificate
chain validated by the PKI of Chapter~24. The Finished messages are
again HMAC. The record layer is AEAD --- AES-GCM or ChaCha20-Poly1305 --- from
Chapter~16, with AES itself from Chapter~12. The random
nonces and ephemeral secrets come from the CSPRNG of Chapter~17. Pull
any one gear and the machine stops; that is why a vulnerability in any single
primitive (a weak RNG, a broken hash, a signature forgery) becomes a TLS break.
\end{connectionbox}

\begin{marginnotebox}
If you want to watch a real handshake, run \texttt{openssl s\_client -connect
example.com:443 -tls1\_3} and read the messages, or capture port 443 in Wireshark
and expand the TLS layer. You will see ClientHello with its key share and SNI in
the clear, ServerHello, then a wall of encrypted handshake records, then
application data. Everything in this chapter is sitting right there on the wire.
\end{marginnotebox}

TLS secures a byte stream once TCP is already moving packets between endpoints, but
what protects the local radio hop before a single TCP packet reaches the router? In
Chapter~26, we drop down to the link layer to see how WEP, WPA2, and WPA3 encrypt and
authenticate WiFi frames broadcast through open air.

\section{Exercises}

\begin{exercise}
\textbf{Trace a tiny handshake's key exchange.} Using $p = 23$, $g = 5$, let the
client pick $a = 4$ and the server pick $b = 9$. Compute the public shares $A$ and
$B$ and the shared secret $Z = g^{ab} \bmod 23$ two ways (as $B^a$ and as $A^b$),
checking they agree. Then explain precisely what Eve sees on the wire and what hard
problem she would have to solve to recover $Z$.
\end{exercise}

\begin{exercise}
\textbf{Why hash the DH output.} The shared secret $Z = g^{ab}$ is an element of
$\Zns{p}$, not a uniform bit string. Give one concrete reason raw $Z$ is unsuitable
as an AES key, and explain what the HKDF-Extract step fixes. (Hint: think about the
distribution of $Z$ and about which subgroup of $\Zns{p}$ it can land in.)
\end{exercise}

\begin{exercise}
\textbf{The role of CertificateVerify.} Suppose a server sent its Certificate but
\emph{not} the CertificateVerify message. Describe a concrete man-in-the-middle
attack an active Eve could mount, and pinpoint which security property
(confidentiality, integrity, or authentication) collapses. Then explain why the
certificate by itself is not enough.
\end{exercise}

\begin{exercise}
\textbf{Finished as downgrade defense.} An attacker edits the ClientHello to remove
all strong cipher suites, leaving only a weak one, hoping the server will pick it.
Trace through the handshake and identify the exact step at which this tampering is
detected. Why does it matter that the Finished MAC is computed over the
\emph{entire} transcript rather than just the final messages?
\end{exercise}

\begin{exercise}
\textbf{Forward secrecy by contrast.} In old RSA-key-transport TLS, the client
encrypts the session secret under the server's long-term RSA public key. Show that
an adversary who records the traffic and \emph{later} obtains the server's RSA
private key can decrypt the session. Then explain, step by step, why the same
adversary fails against ephemeral ECDHE even after stealing the server's signing
key.
\end{exercise}

\begin{exercise}
\textbf{Nonce arithmetic.} A TLS 1.3 connection uses a 12-byte
\texttt{static\_iv} = \texttt{0x01020304\allowbreak 05060708\allowbreak 090A0B0C}. Give the AEAD nonces
for record sequence numbers $0$, $1$, and $258$ (the nonce is
\texttt{static\_iv} $\xor \text{seq}$, with the sequence number right-aligned in
the 12 bytes). Argue that within a single connection no (key, nonce) pair can ever
repeat, and explain why that is essential for AES-GCM specifically.
\end{exercise}

\begin{exercise}
\textbf{Break BEAST on paper.} In CBC mode $C_i = \Enc_k(P_i \xor C_{i-1})$.
Assume Eve knows $C_{i-1}$ (the previous record's last block, used as the next IV
in TLS 1.0) and can choose the next plaintext block $P_i$. She wants to test
whether an unknown secret block $S$ equals a guess $G$. Show how she picks $P_i$ so
that a ciphertext-block collision reveals the answer, and state precisely what
about TLS 1.0's IV choice made this possible. What does TLS 1.3 do that removes the
attack entirely?
\end{exercise}

\begin{exercise}
\textbf{0-RTT replay.} A client uses 0-RTT to send the HTTP request
\texttt{POST /transfer?}\allowbreak\texttt{amount=100\&to=eve} as early data. Describe how an active
attacker exploits the lack of replay protection, what damage results, and the
general rule a server should follow about which requests may be accepted as 0-RTT
data. Contrast this with why the main 1-RTT data is not replayable.
\end{exercise}

\begin{exercise}
\textbf{Heartbleed was not crypto.} Explain why Heartbleed could leak a server's
private key even though every cryptographic algorithm in TLS was implemented
correctly. Identify the class of bug, the role of the unchecked length field, and
the general principle it teaches about the relationship between protocol security
and implementation security.
\end{exercise}

\begin{exercise}
\textbf{Trace the whole machine.} For a single TLS 1.3 connection, list every
cryptographic primitive used and the exact purpose each serves: the key exchange,
the KDF, the signature, the certificate validation, the record cipher, and the
randomness source. For each, name one earlier chapter it came from and state what
breaks in the handshake if that one primitive is broken.
\end{exercise}
